\documentclass[11pt]{article}

\usepackage{amsmath}
\usepackage{amssymb}
\usepackage[utf8]{inputenc}
\usepackage[T1]{fontenc}
\usepackage[margin=1in]{geometry}
\usepackage{hyperref}
\hypersetup{
  colorlinks=true,
  linkcolor=blue!50!black,
  urlcolor=blue!50!black
}

% Code style
\usepackage{listings}
\usepackage{xcolor}
\definecolor{codebg}{RGB}{245,245,247}
\lstdefinestyle{py}{
  basicstyle=\ttfamily\small,
  backgroundcolor=\color{codebg},
  frame=single,
  framerule=0pt,
  rulecolor=\color{codebg},
  keywordstyle=\color{blue!60!black}\bfseries,
  commentstyle=\color{green!40!black}\itshape,
  stringstyle=\color{red!40!black},
  showstringspaces=false,
  breaklines=true,
  language=Python,
  xleftmargin=0.5em,
  xrightmargin=0.5em
}
\lstset{style=py}

% Inline code
\newcommand{\code}[1]{\texttt{\small #1}}

% Vectors and matrices
\newcommand{\vw}{\mathbf{w}}
\newcommand{\vy}{\mathbf{y}}
\newcommand{\vyhat}{\hat{\mathbf{y}}}

\title{HW03 — Linear and Logistic Regression\\\large Specification}
\author{Ramaz ML Course 2026}
\date{}

\begin{document}
\maketitle

\noindent
This document is the \textbf{contract} for HW03. It tells you exactly what to build,
what each method must do, and the shape of the public API your code must match.
The test suite checks this contract directly --- when your code matches the spec,
the tests pass.

\medskip
\noindent
\textbf{How to use this document:} Read it once end-to-end before opening
\code{regression.py}. Then implement methods in roughly the order tests are written
(unit tests first, then \code{fit}, then anything left). Refer back here whenever
you forget a math formula or a return shape.

\tableofcontents

\section{Architecture overview}

Both \code{LinearRegression} and \code{LogisticRegression} follow the same five-step loop:

\begin{center}
\begin{tabular}{ll}
1. & \code{forward(X)} computes predictions $\vyhat$ from inputs $X$. \\
2. & \code{loss(y\_pred, y\_true)} returns a scalar loss. \\
3. & \code{gradient(X, y\_true, y\_pred)} returns $(dw, db)$. \\
4. & Update step: $\vw \leftarrow \vw - \eta \cdot dw$,\quad $b \leftarrow b - \eta \cdot db$. \\
5. & Repeat for \code{epochs} iterations.
\end{tabular}
\end{center}

\noindent
Your \code{fit(X, y, epochs)} method runs this loop \code{epochs} times and returns
the list of loss values, one per epoch. Your \code{predict(X)} method returns the
model's prediction for fresh data:

\begin{itemize}
    \item For linear regression, \code{predict(X)} \emph{is} the forward pass.
    \item For logistic regression, \code{predict(X)} is a hard label ($0$ or $1$)
          obtained by thresholding the forward pass's probability at $0.5$.
\end{itemize}

\section{Standalone helper: \texttt{sigmoid}}

You must implement one free function before either class:

\medskip\noindent
\code{def sigmoid(z: torch.Tensor) -> torch.Tensor}

\medskip\noindent
A non-linear activation that squashes any real number into the
interval $(0, 1)$. \code{LogisticRegression.forward} will apply this to a
linear combination to turn it into a probability.

\medskip\noindent
Element-wise sigmoid:
\[
    \sigma(z) = \frac{1}{1 + e^{-z}}
\]

\noindent
The output is always in the open interval $(0, 1)$, has the same shape as the
input, and satisfies $\sigma(-z) = 1 - \sigma(z)$.

\medskip
\noindent
\code{LogisticRegression.forward} will call this --- but do not call it from
inside \code{LinearRegression.forward}.

\section{LinearRegression}

The required public API is described method-by-method below. Match each
signature exactly --- the tests import the class and call these methods on
instances of it.

\subsection{\texttt{\_\_init\_\_}}

\noindent
\code{def \_\_init\_\_(self, n\_features: int, lr: float = 0.01) -> None}

\medskip\noindent
Set up a fresh model. Called once, when the user does
\code{model = LinearRegression(n\_features=8)}. After this, the model has
parameters initialized to zero but hasn't seen any data yet.

\medskip\noindent
Initialize:
\begin{itemize}
    \item \code{self.w} --- a 1-D tensor of zeros with shape \code{(n\_features,)},
          dtype \code{torch.float32}.
    \item \code{self.b} --- a scalar tensor (0-D) with value $0.0$, dtype
          \code{torch.float32}.
    \item \code{self.lr} --- store the learning rate as-is.
\end{itemize}

The grader reads \code{self.w} and \code{self.b} after \code{fit}, so use these
exact attribute names.

\subsection{\texttt{forward}}

\noindent
\code{def forward(self, X: torch.Tensor) -> torch.Tensor}

\medskip\noindent
Step 1 of every training iteration. Take the current parameters
and produce a prediction for each row of $X$. Also called at inference time
(via \code{predict}) on new data.

\medskip\noindent
Compute the linear prediction:
\[
    \vyhat = X \vw + b
\]

\begin{itemize}
    \item Input shape: \code{X} has shape \code{(n, d)} where
          \code{d == self.w.shape[0]}.
    \item Output shape: \code{(n,)}.
\end{itemize}

\subsection{\texttt{loss}}

\noindent
\code{def loss(self, y\_pred: torch.Tensor, y\_true: torch.Tensor) -> torch.Tensor}

\medskip\noindent
Step 2 of every training iteration. Collapse the per-example
prediction error into one scalar. The training loop tries to make this number
smaller.

\medskip\noindent
Mean squared error:
\[
    \mathrm{MSE} = \frac{1}{n} \sum_{i=1}^{n} (\hat{y}_i - y_i)^2
\]

\begin{itemize}
    \item Inputs: both shape \code{(n,)}.
    \item Output: a 0-dimensional (scalar) tensor.
\end{itemize}

\subsection{\texttt{gradient}}

\noindent
\code{def gradient(self, X: torch.Tensor, y\_true: torch.Tensor, y\_pred: torch.Tensor)\\
\hspace*{2em}-> tuple[torch.Tensor, torch.Tensor]}

\medskip\noindent
Step 3 of every training iteration. Figure out which direction
to move each parameter to reduce the loss. The training loop will subtract a
small multiple (\code{lr}) of these gradients from $\vw$ and $b$.

\medskip\noindent
Returns the gradients of the MSE loss with respect to $\vw$ and $b$. Use the
formula derived in lecture (no \code{torch.autograd}, no \code{.backward()}):
\[
    \frac{\partial L}{\partial \vw} = \frac{2}{n}\, X^{\top} (\vyhat - \vy)
    \qquad
    \frac{\partial L}{\partial b} = \frac{2}{n} \sum_{i=1}^{n} (\hat{y}_i - y_i)
\]

\begin{itemize}
    \item \code{dw} has shape \code{(d,)}.
    \item \code{db} is a scalar tensor.
\end{itemize}

\subsection{\texttt{fit}}

\noindent
\code{def fit(self, X: torch.Tensor, y: torch.Tensor, epochs: int = 200) -> list[float]}

\medskip\noindent
Run the full training loop. This is the only method a normal
user will call directly --- it orchestrates \code{forward}, \code{loss},
\code{gradient}, and the parameter update step in sequence.

\medskip\noindent
Full-batch gradient descent. The algorithm:

\begin{enumerate}
    \item Initialize an empty history list.
    \item Repeat for \code{epochs} iterations:
    \begin{enumerate}
        \item Compute predictions from the current parameters (call your own
              \code{forward}).
        \item Compute the loss between predictions and targets (call your own
              \code{loss}).
        \item Append the loss \emph{as a Python float} (not a tensor) to the
              history.
        \item Compute the gradients (call your own \code{gradient}).
        \item Update parameters by stepping against the gradient:
              $\vw \leftarrow \vw - \eta \cdot dw$,\quad
              $b \leftarrow b - \eta \cdot db$.
    \end{enumerate}
    \item Return the history list.
\end{enumerate}

\begin{itemize}
    \item Returns a \code{list[float]} of length exactly \code{epochs}.
    \item After \code{fit} returns, \code{self.w} and \code{self.b} hold the
          trained parameters.
\end{itemize}

\subsection{\texttt{predict}}

\noindent
\code{def predict(self, X: torch.Tensor) -> torch.Tensor}

\medskip\noindent
Inference. Produce predictions on new data using the trained
parameters. Tests call this after \code{fit} to evaluate model quality.

\medskip\noindent
For linear regression, \code{predict(X)} is identical to \code{forward(X)}.
(It exists so that \code{LinearRegression} and \code{LogisticRegression} share
a uniform inference API.)

\subsection*{Invariants the tests will check}

\begin{itemize}
    \item For sufficiently small \code{lr}, \code{loss\_history} is non-increasing.
    \item On data generated by $y = X w_{\text{true}} + b_{\text{true}} + \varepsilon$
          with small $\varepsilon$, \code{fit} for enough epochs recovers
          $w_{\text{true}}$ and $b_{\text{true}}$ to within tolerance.
    \item \code{forward(X)} returns shape \code{(n,)}; \code{gradient} returns
          \code{(dw, db)} with shapes \code{(d,)} and scalar; \code{predict(X)}
          returns shape \code{(n,)}.
\end{itemize}

\section{LogisticRegression}

The required public API is described method-by-method below. Match each
signature exactly.

\subsection{\texttt{\_\_init\_\_}}

\noindent
\code{def \_\_init\_\_(self, n\_features: int, lr: float = 0.5) -> None}

\medskip\noindent
Same as \code{LinearRegression.\_\_init\_\_} --- set up a fresh model.

\medskip\noindent
Same behavior as \code{LinearRegression.\_\_init\_\_}. Default \code{lr} is
larger because the BCE gradient is on a smaller scale --- see Question 3 in
the writeup.

\subsection{\texttt{forward}}

\noindent
\code{def forward(self, X: torch.Tensor) -> torch.Tensor}

\medskip\noindent
Step 1 of every training iteration. Produce a class-1
probability for each input row.

\medskip\noindent
Probability of class~1:
\[
    \vyhat = \sigma(X \vw + b)
\]

Use the standalone \code{sigmoid} you implemented in §2.

\begin{itemize}
    \item Output shape: \code{(n,)}. All values lie strictly in $(0, 1)$.
\end{itemize}

\subsection{\texttt{loss}}

\noindent
\code{def loss(self, y\_pred: torch.Tensor, y\_true: torch.Tensor) -> torch.Tensor}

\medskip\noindent
Step 2 of every training iteration. Measure prediction error for
a binary classification problem. BCE penalizes confident wrong predictions
much more than uncertain ones.

\medskip\noindent
Binary cross-entropy:
\[
    \mathrm{BCE}
    = -\frac{1}{n} \sum_{i=1}^{n}
        \big[
            y_i \log \hat{y}_i + (1 - y_i) \log (1 - \hat{y}_i)
        \big]
\]

\begin{itemize}
    \item \code{y\_true} values are in $\{0, 1\}$ (encoded as floats).
    \item \textbf{Numerical safety:} clamp \code{y\_pred} into
          $[\varepsilon, 1 - \varepsilon]$ with $\varepsilon = 10^{-7}$ before
          taking \code{log}. Otherwise $\log(0) = -\infty$ will blow up your loss.
\end{itemize}

\subsection{\texttt{gradient}}

\noindent
\code{def gradient(self, X: torch.Tensor, y\_true: torch.Tensor, y\_pred: torch.Tensor)\\
\hspace*{2em}-> tuple[torch.Tensor, torch.Tensor]}

\medskip\noindent
Step 3 of every training iteration. Same purpose as the linear
case --- decide which direction to move each parameter to reduce the loss.

\medskip\noindent
Same shape as the linear case, with the factor changed:
\[
    \frac{\partial L}{\partial \vw} = \frac{1}{n}\, X^{\top} (\vyhat - \vy)
    \qquad
    \frac{\partial L}{\partial b} = \frac{1}{n} \sum_{i=1}^{n} (\hat{y}_i - y_i)
\]

\noindent
Why no factor of 2 here? See writeup Question 3 --- the sigmoid derivative
cancels something in the BCE derivative. Lesson 4 has the full derivation.

\subsection{\texttt{fit}}

\noindent
\code{def fit(self, X: torch.Tensor, y: torch.Tensor, epochs: int = 300) -> list[float]}

\medskip\noindent
Run the full training loop --- same orchestration as
\code{LinearRegression.fit}.

\medskip\noindent
Same training loop as \code{LinearRegression.fit}. The defaults
(\code{lr=0.5}, \code{epochs=300}) are tuned for the Iris classification dataset
used in the analysis.

\subsection{\texttt{predict\_proba}}

\noindent
\code{def predict\_proba(self, X: torch.Tensor) -> torch.Tensor}

\medskip\noindent
Inference, soft version. Returns the probability of class 1
without thresholding. Useful when callers want a confidence value.

\medskip\noindent
Returns the same thing as \code{forward(X)} --- a vector of probabilities in
$(0, 1)$. This method exists so callers can read soft predictions without the
threshold being applied.

\subsection{\texttt{predict}}

\noindent
\code{def predict(self, X: torch.Tensor) -> torch.Tensor}

\medskip\noindent
Inference, hard version. Returns the predicted class label
($0$ or $1$) for each input.

\medskip\noindent
Apply the standard $0.5$ threshold:
\[
    y_i =
    \begin{cases}
        1 & \text{if } \hat{y}_i > 0.5 \\
        0 & \text{otherwise}
    \end{cases}
\]

\begin{itemize}
    \item Output shape: \code{(n,)}, values in $\{0, 1\}$.
    \item Output dtype should match \code{y\_true.dtype} (\code{float32} in our
          datasets).
\end{itemize}

\subsection*{Invariants the tests will check}

\begin{itemize}
    \item \code{forward} outputs are strictly in $(0, 1)$.
    \item For sufficiently small \code{lr}, \code{loss\_history} is non-increasing.
    \item On linearly separable data, \code{predict} achieves $>95\%$ accuracy
          after enough epochs.
    \item \code{predict(X)} agrees with \code{(predict\_proba(X) > 0.5).float()}.
\end{itemize}

\section{Analysis (\texttt{analysis.py})}

\code{analysis.py} contains four utility functions used by the test suite and
by your final \code{analysis.py} run that produces plots for \code{writeup.md}.
These are utility helpers --- the learning goal is the regression classes
above, not data preprocessing --- so the spec names the sklearn entry points
directly. Match the signatures exactly.

\subsection{\texttt{load\_regression\_data}}

\noindent
\code{def load\_regression\_data() -> tuple[Tensor, Tensor, Tensor, Tensor]}

\medskip\noindent
Load the California Housing dataset and return train/test splits as float32
tensors.

\medskip\noindent
Steps:
\begin{enumerate}
    \item Call \code{sklearn.datasets.fetch\_california\_housing()}.
    \item Split with \code{sklearn.model\_selection.train\_test\_split(test\_size=0.2,
          random\_state=42)}.
    \item Standardize features: fit a \code{StandardScaler} on the training set
          only, then transform both sets.
    \item Convert all four arrays to \code{torch.Tensor} with
          \code{dtype=torch.float32}.
\end{enumerate}

\noindent
Return order: \code{(X\_train, X\_test, y\_train, y\_test)}. Shapes:
\code{X\_train} is \code{(n\_train, 8)}, \code{X\_test} is \code{(n\_test, 8)},
labels are 1-D.

\subsection{\texttt{train\_linear\_model}}

\noindent
\code{def train\_linear\_model(X\_train: Tensor, y\_train: Tensor,\\
\hspace*{2em}lr: float = 0.01, epochs: int = 200) -> tuple[LinearRegression, list[float]]}

\medskip\noindent
Thin wrapper: construct a fresh \code{LinearRegression} with
\code{n\_features=X\_train.shape[1]} and the given \code{lr}, call its
\code{fit} method for \code{epochs} epochs, and return
\code{(model, loss\_history)}.

\subsection{\texttt{load\_classification\_data}}

\noindent
\code{def load\_classification\_data() -> tuple[Tensor, Tensor, Tensor, Tensor]}

\medskip\noindent
Load a 2-class subset of the Iris dataset (setosa vs.\ versicolor), keeping
only the first two features so the decision boundary can be plotted in 2D.

\medskip\noindent
Steps:
\begin{enumerate}
    \item Call \code{sklearn.datasets.load\_iris()}.
    \item Filter to rows where \code{target < 2}; keep only the first two
          feature columns.
    \item Split with \code{train\_test\_split(test\_size=0.2, random\_state=42)}.
    \item Standardize features (fit on train only).
    \item Convert all four arrays to \code{float32} torch tensors.
\end{enumerate}

\noindent
Labels are in $\{0, 1\}$ encoded as floats.

\subsection{\texttt{train\_logistic\_model}}

\noindent
\code{def train\_logistic\_model(X\_train: Tensor, y\_train: Tensor,\\
\hspace*{2em}lr: float = 0.5, epochs: int = 300) -> tuple[LogisticRegression, list[float]]}

\medskip\noindent
Same shape as \code{train\_linear\_model} --- construct a fresh
\code{LogisticRegression}, call \code{fit}, return \code{(model, loss\_history)}.

\section{How to read the test file}

\code{test\_regression.py} is organized in three layers --- work on them in order:

\subsection*{\texttt{@pytest.mark.unit} --- one method at a time}

Each test class targets one piece of one class. Get these green first; they
verify the math primitives independently (forward shapes, loss values,
gradients matching the formula) without requiring a working training loop.

\subsection*{\texttt{@pytest.mark.integration} --- \texttt{fit} end-to-end}

Once unit tests pass, move on. Integration tests run a full \code{fit} on
synthetic data and check that the loss decreases, that the recovered parameters
are close to ground truth, and that \code{predict} produces sensible outputs.

\subsection*{\texttt{@pytest.mark.analysis} --- \texttt{analysis.py} helpers}

Tests for the four helper functions in \code{analysis.py} (data loading and
training wrappers).

\medskip
\noindent
Run a single layer with:

\begin{lstlisting}[language=bash]
uv run pytest -m unit         # just unit tests
uv run pytest -m integration  # just integration
uv run pytest -m analysis     # just analysis helpers
\end{lstlisting}

\section{Restrictions}

\begin{itemize}
    \item \textbf{No \code{torch.autograd}, no \code{.backward()}, no
          \code{requires\_grad=True} anywhere in \code{regression.py}.} You are
          computing every gradient by hand from the derivations in lecture. The
          \code{conftest.py} fails the run if it detects autograd usage in your file.
    \item \textbf{No \code{sklearn.linear\_model}.} Same reason --- implementing
          logistic regression yourself is the point of the assignment.
\end{itemize}

\end{document}
